\clearpage
\section{Appendix: Liturgical Resolution}
The celebration is the Twenty-eighth Sunday in Ordinary Time, Year A, in the postconciliar Roman Rite for the dioceses of the United States of America, on 11 October 2026. It is a Sunday of Ordinary Time; the color is green. The governing books are the U.S. \work{Roman Missal}, Third Edition (2011), with the emended third typical Latin edition (2008), and the U.S. \work{Lectionary for Mass}, second typical edition, Volume I, no. 142. The 1981 second typical \work{Ordo lectionum Missae}, no. 142, and the official dated U.S. readings agree on the appointments studied here.

The permanent parent is PC-S54 and the formula is PC-S54-A. Its common Missal formulary belongs to Ordinary Time Week 28 and serves all three Sunday cycles. The national calendar establishes this dated Sunday; no particular diocesan, parish, titular, dedication or religious-institute overlay was supplied. Weekday Cycle II is independently applicable to the year's ferial readings and supplies none of this Sunday's readings. The \work{General Introduction to the Lectionary}, nos. 66--68, 106--107, distinguishes the Old Testament--Gospel correlation from the semi-continuous apostolic course.

\begin{description}[style=nextline,leftmargin=0pt,labelsep=0pt]
\item[\bid{gospel-long-or-short}] Appointed alternative: Mt 22:1--14 or 1--10 (Lectionary 142; GIRM 360). Both are studied; no local selection is claimed. The short form omits the garment, expulsion and concluding saying.
\item[\bid{communion-psalm-or-john}] Appointed alternative: cf. Ps 33:11 or 1 Jn 3:2 in the Missal. Both are studied; they are mutually alternative, and neither is a documented local choice.
\item[\bid{entrance-missal-or-permitted-chant}] GIRM 48 governs permitted chant routes. The Missal antiphon is the studied route; its recitation applies when there is no singing. No alternative chant is selected.
\item[\bid{psalm-lectionary-or-permitted-sung-form}] GIRM 61 permits defined sung options. The no. 142 stanzas and response are studied; another form is unselected.
\item[\bid{acclamation-lectionary-or-graduale}] GIRM 62 governs the source routes. The no. 142 Alleluia verse is studied. The one-reading circumstance in GIRM 63(c) is not this three-reading Sunday.
\item[\bid{sprinkling-in-place-of-penitential-act}] Permitted under GIRM 51 and Missal Appendix II; unselected. No additional sprinkling texts are included.
\item[\bid{preparation-chant}] Permitted under GIRM 74, with 48; unselected. No proper Missal Offertory antiphon or concrete local song is identified.
\item[\bid{preface-eucharistic-prayer}] Lawful coupled choices remain open under GIRM 364--365. No unique Week 28 Preface is appointed. Eucharistic Prayer IV, if chosen lawfully, retains its invariable Preface.
\item[\bid{communion-missal-or-permitted-chant}] GIRM 87--88 governs the permitted chant and recitation routes; no additional local chant or post-Communion song is selected.
\item[\bid{concluding-blessing}] The ordinary conclusion applies (GIRM 90); no concrete solemn blessing or prayer over the people is selected.
\end{description}
Gloria, Creed, homily and Universal Prayer retain their ordinary Sunday places; they do not become new proper appointments. No Sequence, proper-specific Eucharistic Prayer insertion or distinctive conclusion is established. The complete branches and shared-owner relationship are recorded in the instance and source audits.

\section{Appendix: Scope and Qualifications}
This expansive study treats the full national-calendar Year A appointment, including both Gospel lengths and both Communion alternatives. Its two interpretations are editorial readings of the assembled texts, supported by the particular arguments attributed to their witnesses. No Father is presented as the author of an interpretation of this modern whole formulary. The interpretations' four senses concern Scripture directly and extend to the composed prayers by theological analogy. The bounded source inquiry is recorded in the canonical research scope and interpretation audit; it is not a claim to exhaust all patristic or saintly reception.

The biblical study text is public-domain Douay--Rheims/Challoner English, Project Gutenberg 1581 witness. It is not the approved proclaimed English: the U.S. Lectionary controls the celebration. Missal and Lectionary translations are not reproduced or reconstructed here. Minimal Latin incipits identify the orations. Their subjects are original explanations, not substitute English prayers. The modern Corpus Thomisticum deliveries supply attributed analysis of Thomas's Latin, not a claim to republish those deliveries as public-domain text.

Common Missal evidence is held by the canonical Week 28 owner. The inspected Latin source is a secondary digital reproduction of the 2002 third typical edition. The official ICEL Antiphonary identifies the antiphon layer; the FDLC compilation supplies approved excerpts of the Collect and Prayer after Communion. The inspected 2008 \work{Notitiae} variation list records no Week XXVIII change in its Ordinary Time block, but is selective rather than an exhaustive collation. Exact publisher-specific 2008/2011 altar-book collation and an exact U.S. English Offerings exemplar remain unavailable. FDLC's Collect conclusion predates the current U.S. wording; its complete text is not offered for recitation. Claims of ancient prayer derivation remain uncollated and are not asserted.

Gregory and Thomas were consulted in the specified Latin digital witnesses; Augustine, Chrysostom, Bellarmine and the mystagogical lecture in identified historical English. The NPNF psalm exposition and O'Sullivan's Bellarmine are abridged witnesses. No fresh Greek critical collation or collation of those translations against the Latin is claimed. NPNF's Matthew homily note 2585 qualifies the biblical lemma's printed extent; the subsequent garment exposition is present in the homily. The proposed royal distribution of wedding garments was not established by the inspected evidence. A later direct saintly exposition of the precise Johannine Communion verse was not established in the bounded sweep; Augustine supplies its direct reception, and Thomas's banquet account a compatible theological relation.

The historical chronology sources preserve alternatives, uncertain places and distinct subjects. Their ranges do not claim contemporary scholarly unanimity. The study's present source boundary is 6 October 2026. This is a Catholic study aid, not an official liturgical book, an ecclesiastically approved commentary or a substitute for competent liturgical direction. Source inspection and author proof checks do not constitute ecclesiastical approval.

\section{References}
\begin{itemize}
\item \work{Holy Bible}, Douay--Rheims, Challoner revision, Project Gutenberg eBook 1581; Psalms 22, 33, 129; Isaiah 24--26; Matthew 21--22; Ephesians 1; Philippians 4; 1 John 2--3. Public-domain study witness.
\item Augustine, \work{Enarrationes in Psalmos} 22.1--6; 33, second exposition, 8--14; 129.1--8. NPNF1 VIII, Hebrew-numbered headings XXIII, XXXIV and CXXX; historical CCEL English witness. Also \work{Homilies on the First Epistle of John} IV.4--9, NPNF1 VII, historical English CCEL witness.
\item Gregory the Great, \work{Homiliae in Evangelia} 38.1--16, especially 1, 5, 7--14. Latin Wikisource witness, complete EPUB and text extraction acquired 17 September 2026; paraphrased here. The export's attribution and CC BY-SA 3.0 terms govern its modern transcription and encoding separately from the ancient Latin.
\item John Chrysostom, \work{Homilies on Matthew} LXIX.1--3, NPNF1 X; \work{Homilies on Philippians} XV, ad 4:10--23, and \work{Homilies on Ephesians} III, ad 1:15--22, NPNF1 XIII. Historical English CCEL witnesses.
\item Thomas Aquinas, \work{Super Psalmo} 22, nn. 1--2, and 33, nn. 10--11; \work{Expositio super Isaiam ad litteram}, cap. 25 (paragraph 84985); \work{Super Philippenses} 4, lectio 2 (87844); \work{Super Ephesios} 1, lectio 6 (87792). Latin Corpus Thomisticum witnesses; original summaries of the specified loci.
\item Robert Bellarmine, \work{A Commentary on the Book of Psalms}, Psalm 129, trans. John O'Sullivan (Dublin, 1866), abridged English; retained historical-text witness from e-Catholic 2000.
\item \work{Mystagogical Catechesis} III.1--7, transmitted under Cyril of Jerusalem, numbered \work{Catechetical Lecture} XXI, NPNF2 VII; historical English CCEL witness.
\item \work{Ordo lectionum Missae}, editio typica altera (1981), no. 142, p. 77; \work{General Introduction to the Lectionary}, nos. 65--69, 74--77, 103--107, approved English hosted by the Liturgy Office of England and Wales.
\item USCCB, \work{Liturgical Calendar for the Dioceses of the United States of America 2026}, pp. 5, 40; \href{https://bible.usccb.org/bible/readings/101126.cfm}{readings for 11 October 2026}, no. 142. Official dated witnesses inspected 6 October 2026.
\item \work{General Instruction of the Roman Missal}, U.S. English, chapters II and VII, nos. 48, 51, 53--54, 61--71, 74, 87--90, 352--367; official USCCB states inspected 6 October 2026.
\item \work{Missale Romanum}, third typical edition (2002), secondary digital reproduction, physical PDF p. 291; ICEL \work{Antiphonary} (2010), official CBCEW excerpt, p. 80; FDLC, \work{Mystagogical Reflections}, Ordinary Time 24--30 (2016), pp. 14--15; \work{Notitiae} 44, nos. 503--504 (2008), pp. 367, 371. Bounded common-formulary witnesses.
\item \work{Catholic Encyclopedia}: Charles Souvay, ``Isaias'' (1910); Ferdinand Prat, ``St. Paul'' (1911); Achille Vander Heeren, ``Epistle to the Philippians'' (1911); Paulin Ladeuze, ``Epistle to the Ephesians'' (1909); Anthony Maas, ``Chronology of the Life of Jesus Christ'' (1910); E. Jacquier, ``Gospel of St. Matthew'' (1911); Alfred Durand, ``The New Testament'' (1912); Walter Drum, ``Epistles of Saint John'' (1910); John Corbett, ``David, King'' (1908). Chronological discussions at the loci named in the historical appendix and canonical chronology record.
\item NABRE introductions to Psalms and Matthew, official USCCB witnesses; historical setting, composition boundaries and the distinct critical Matthew comparison. Biblical translation not reproduced.
\item George Leo Haydock, annotated Douay Bible, inspected facsimile witness, physical pp. 396, 398--399, annotations and year headings around 1 Kings 20--22; narrative frame for the Psalm 33 superscription setting.
\end{itemize}
